\documentclass[10pt]{article}
\usepackage[T1]{fontenc}
\usepackage[utf8]{inputenc}
\usepackage[french]{babel}
\usepackage{amsmath,amssymb,mathrsfs,mathpazo}
\usepackage{geometry}
\geometry{a4paper,margin=2cm}
\usepackage{graphicx,array,tabularx,tabularray,colortbl,multirow,multicol,booktabs,enumitem,pifont,ulem,fancybox,diagbox,cancel}
\usepackage{pstricks,pstricks-add,pst-plot,pst-tree,pst-node,pst-text,pst-func,pst-math,pst-bspline,pst-3dplot,pst-eucl,pst-solides3d}
\let\pscalculate\undefined
\usepackage{pst-bezier}
\usepackage{tikz,pgfplots,tkz-tab}
\usetikzlibrary{arrows,arrows.meta,patterns,calc,positioning,trees,decorations.markings,shapes.geometric,shapes.misc,shapes.arrows,intersections,decorations.pathmorphing,decorations.pathreplacing,backgrounds,shadows,babel}
\pgfplotsset{compat=1.18}
\usepgfplotslibrary{fillbetween}
\usepackage[np]{numprint}
\usepackage{listings,makecell,pas-tableur,twoopt}
\let\barre\undefined
\usepackage[pstricks]{tablvar}
\definecolor{darkgreen}{RGB}{0,100,0}
\providecommand{\red}{\color{red}}
\providecommand{\blue}{\color{blue}}
\providecommand{\green}{\color{green}}
\providecommand{\black}{\color{black}}
\providecommand{\white}{\color{white}}
\providecommand{\gray}{\color{gray}}
\providecommand{\cyan}{\color{cyan}}
\providecommand{\magenta}{\color{magenta}}
\providecommand{\yellow}{\color{yellow}}
\providecommand{\violet}{\color{violet}}
\providecommand{\orange}{\color{orange}}
\providecommand{\euro}{\texteuro}
\providecommand{\EUR}{\texteuro}
\providecommand{\ieme}{\textsuperscript{e}}
\providecommand{\degres}{\ensuremath{^\circ}}
\providecommand{\degree}{\ensuremath{^\circ}}
\providecommand{\vv}[1]{\overrightarrow{#1}}
\providecommand{\Coord}[1]{#1}
\providecommand{\decosix}{}
\providecommand{\starredbullet}{\ensuremath{\bullet}}
\pagestyle{empty}
\setlength{\parindent}{0pt}
\usepackage{adjustbox}\DeclareUnicodeCharacter{2236}{:}\DeclareUnicodeCharacter{2190}{\ensuremath{\leftarrow}}
\providecommand{\R}{\mathbb{R}}
\providecommand{\N}{\mathbb{N}}
\providecommand{\Z}{\mathbb{Z}}
\providecommand{\Q}{\mathbb{Q}}
\providecommand{\D}{\mathbb{D}}
\providecommand{\C}{\mathbb{C}}
\providecommand{\e}{{\mathrm{e}}}
\providecommand{\vect}[1]{\overrightarrow{#1}}
\providecommand{\vectt}[1]{\overrightarrow{\mathrm{#1}}}
\providecommand{\vvt}[1]{\overrightarrow{\mathrm{#1}}}
\providecommand{\np}[2][]{{#2~#1}}
\providecommand{\vv}[1]{{\overrightarrow{#1}}}
\renewcommand{\d}{{\mathrm{d}}}
\providecommand{\degres}{{^\circ}}
\providecommand{\degre}{{^\circ}}
\providecommand{\blue}{\color{blue}}
\providecommand{\red}{\color{red}}
\providecommand{\enskip}{\quad}
\providecommand{\llbracket}{\lbrack\!\lbrack}
\providecommand{\rrbracket}{\rbrack\!\rbrack}
\providecommand{\euro}{{\text{€}}}
\providecommand{\no}{{\text{n°}}}
\providecommand{\interval}[2]{{[#1;#2]}}
\providecommand{\intervFO}[2]{{[#1;#2[}}
\providecommand{\dotfill}{\ldots}

\providecommand{\ds}{\displaystyle}
\providecommand{\textding}[1]{#1}
\providecommand{\Oij}{\(\mathrm{O};\vec{\imath},\vec{\jmath}\)}
\providecommand{\Oijk}{\(\mathrm{O};\vec{\imath},\vec{\jmath},\vec{k}\)}
\providecommand{\Ouv}{\(\mathrm{O};\vec{u},\vec{v}\)}
\providecommand{\intervFF}[2]{\ensuremath{[#1;#2]}}
\renewcommand{\euro}{\ifmmode\text{\texteuro}\else\texteuro\fi}
\DeclareUnicodeCharacter{00B0}{\ifmmode^{\circ}\else\textdegree\fi}
\DeclareUnicodeCharacter{20AC}{\ifmmode\text{\texteuro}\else\texteuro\fi}
% Énoncés : Éducation nationale. Transcription et corrigés : APMEP, Denis Vergès et contributeurs. Découpage et classement : CoopMaths, Sébastien Lozano, Éric Elter et contributeurs. Adaptation de la présentation : Maths Explorer.
% CC BY-SA : https://coopmaths.fr/www/about/
% https://coopmaths.fr/alea/static/bac/2024/tex/bac_2024_11_sujet1_amsud_2.tex
\begin{document}
\noindent\textit{APMEP / CoopMaths. CC BY-SA. Présentation adaptée par Maths Explorer.}\par\medskip
\medskip

On dispose de deux urnes opaques U$_1$ et U$_2$.

L'urne U$_1$ contient 4 boules noires et 6 boules blanches.

L'urne U$_2$ contient 1 boule noire et 3 boules blanches.

On considère l'expérience aléatoire suivante :

On pioche au hasard une boule dans U$_1$ que l'on place dans U$_2$, puis on pioche au hasard une boule dans U$_2$.

On note:

\begin{itemize}[label=$\bullet~$]
\item $N_1$ l'évènement \og Piocher une boule noire dans l'urne U$_1$ \fg.
\item $N_2$ l'évènement \og Piocher une boule noire dans l'urne U$_2$ \fg.
\end{itemize}

Pour tout évènement $A$, on note $\overline{A}$ son évènement contraire.

\medskip

\textbf{PARTIE A}

\medskip

\begin{minipage}{0.53\linewidth}

\begin{enumerate}
\item On considère l'arbre de probabilités ci-contre.
	\begin{enumerate}
		\item Justifier que la probabilité de piocher une boule noire dans l'urne U$_2$ sachant qu'on a pioché une boule blanche dans l'urne U$_1$ est $0,2$.
		\item Recopier et compléter l'arbre de probabilités ci-contre, en faisant apparaître sur chaque branche les probabilités des évènements concernés, sous forme décimale.
	\end{enumerate}
\end{enumerate}
\end{minipage}\hfill
\begin{minipage}{0.43\linewidth}
\begin{center}
\pstree[treemode=R,nodesepA=0pt,nodesepB=2.5pt,treesep = 1cm,levelsep=2.5cm]{\TR{}}
{\pstree{\TR{$N_1$~}}
	{\TR{$N_2$}
	\TR{$\overline{N_2}$}	
	}
\pstree{\TR{$\overline{N_1}$~}}
	{\TR{$N_2$}\ncput*{0,2}
	\TR{$\overline{N_2}$}	
	}
}
\end{center}
\end{minipage}

\begin{enumerate}[start=2]
\item Calculer la probabilité de piocher une boule noire dans l'urne U$_1$ et une boule noire dans l'urne U$_2$.
\item Justifier que la probabilité de piocher une boule noire dans l'urne U$_2$ est égale à 0,28.
\item On a pioché une boule noire dans l'urne U$_2$

Calculer la probabilité d'avoir pioché une boule
blanche dans l'urne U$_1$ . On donnera le résultat sous forme décimale arrondie à $10^{-2}$.
\end{enumerate}

\medskip

\textbf{PARTIE B}

\medskip

$n$ désigne un entier naturel non nul.

L'expérience aléatoire précédente est répétée $n$ fois de façon identique et indépendante, c'est-à-dire que les urnes U$_1$ et U$_2$ sont remises dans leur configuration initiale, avec respectivement 4 boules noires et 6 boules blanches dans l'urne U$_1$ et 1 boule noire et 3 boules blanches dans l'urne U$_2$, entre chaque expérience.

On note $X$ la variable aléatoire qui compte le nombre de fois où on pioche une boule noire dans l'urne U$_2$.

On rappelle que la probabilité de piocher une boule noire dans l'urne U$_2$ est égale à 0,28 et celle de piocher une boule blanche dans l'urne U$_2$ est égale à 0,72.

\medskip

\begin{enumerate}
\item Déterminer la loi de probabilité suivie par $X$. Justifier votre réponse.
\item Déterminer par le calcul le plus petit entier naturel n tel que: 
\[1 - 0,72^n \geqslant 0,9.\]
\item Interpréter le résultat précédent dans le contexte de l'expérience.
\end{enumerate}

\medskip

\textbf{PARTIE C}

\medskip

Dans cette partie les urnes U$_1$ et U$_2$ sont remises dans leur configuration initiale, avec respectivement 4 boules noires et 6 boules blanches dans l'urne U$_1$ et 1 boule noire et 3 boules blanches dans l'urne U$_2$.

\medskip

On considère la nouvelle expérience aléatoire suivante :

On pioche simultanément deux boules dans l'urne U$_1$ que l'on place dans l'urne U$_2$, puis on pioche au hasard une boule dans l'urne U$_2$.

\medskip

\begin{enumerate}
\item Combien y a-t-il de tirages possibles de deux boules simultanément dans l'urne U$_1$ ?
\item Combien y a-t-il de tirages possibles de deux boules simultanément dans l'urne U$_1$
contenant exactement une boule blanche et une boule noire ?
\item La probabilité de piocher une boule noire dans l'urne U$_2$ avec cette nouvelle expérience est- elle supérieure à la probabilité de tirer une boule noire dans l'urne U$_2$ avec l'expérience de la partie A ? Justifier votre réponse.

\emph{On pourra s'aider d'un arbre pondéré modélisant cette expérience}.
\end{enumerate}

\bigskip

\end{document}
